\chapter{Transformasi Fourier}\label{ch:b3:fouriertransform}

Deret Fourier menguraikan sinyal periodik menjadi harmonik
diskret; sedangkan \emph{transformasi} Fourier melakukan hal yang sama bagi
sinyal pada seluruh garis, dengan frekuensi yang \omterm{def:b3:topology:continuity}{kontinu}. Ia
mengubah pendiferensialan menjadi perkalian, konvolusi menjadi
hasil kali, dan fungsi Gauss menjadi fungsi Gauss --- yakni alasan ia
menyelesaikan persamaan diferensial, menggerakkan pengolahan sinyal, dan
akan membuktikan teorema limit pusat di \cref{ch:b3:clt}. Bab
ini mengembangkan teori $L^1$-nya (Riemann--Lebesgue,
pembalikan, keinjektifan), \omterm{def:b3:fouriertransform:schwartz}{kelas Schwartz} yang
transformasinya menjadi bijeksi \omterm{prop:b3:galois:perfect}{sempurna}, dan teori $L^2$-nya
(Plancherel: bahwa transformasinya, sampai sebuah konstanta, merupakan operator
uniter), beserta dua terapan andalan: yakni persamaan panas,
yang diselesaikan sampai tuntas pada soal akhir pekan, dan rumus penjumlahan
Poisson. Kesepakatannya:
\[
\hat f(\xi) = \int_\R f(x)\,\eu^{-\iu\xi x}\,\dd x .
\]

\section{Transformasinya pada \texorpdfstring{$L^1$}{L1}}

\begin{proposition}\label{prop:b3:fouriertransform:basic}
Untuk $f \in L^1(\R)$: $\hat f$ terdefinisi dengan baik, terbatas
(sebab $\norm{\hat f}_\infty \leq \norm f_1$), \omterm{def:b3:topology:continuity}{kontinu}, dan:
\begin{enumerate}
  \item $\widehat{\tau_af}(\xi) = \eu^{-\iu a\xi}\hat f(\xi)$
        dan $\widehat{\eu^{\iu ax}f}(\xi) = \hat f(\xi - a)$;
  \item $\widehat{f(\cdot/\lambda)}(\xi) = \lambda\hat
        f(\lambda\xi)$ untuk $\lambda > 0$;
  \item jika $xf \in L^1$, maka $\hat f$ bersifat $\mathcal C^1$ dengan
        $(\hat f)'(\xi) = \widehat{(-\iu x)f}(\xi)$;
  \item jika $f \in \mathcal C^1$ dengan $f' \in L^1$ (dan $f \to
        0$ di $\pm\infty$, yang di sini otomatis), maka
        $\widehat{f'}(\xi) = \iu\xi\hat f(\xi)$;
  \item $\widehat{f * g} = \hat f\,\hat g$ untuk $f, g \in L^1$.
\end{enumerate}
\end{proposition}

\begin{proof}
Keterbatasannya: $\abs{\hat f} \leq \int\abs f$. \omterm{def:b3:topology:continuity}{Kekontinuannya}: lewat kekonvergenan
terdominasi dengan pendominasi $\abs f$ (\cref{thm:b3:lebesgue:paramcont}).
(1), (2): lewat substitusinya
(\cref{thm:b3:product:linearchange}). (3): lewat pendiferensialan
di bawah integralnya, dengan pendominasi $\abs{xf}$
(\cref{thm:b3:lebesgue:paramdiff}). (4): mula-mula, $f(x) = f(0)
+ \int_0^xf'$ berlimit di $\pm\infty$ (sebab $f' \in L^1$), yang
haruslah $0$ (sebab $f \in L^1$); lalu integralkan secara parsial pada $[-A,
A]$ dan biarkan $A \to \infty$. (5): lewat Fubini, yang sah sebab
$(x,y)\mapsto f(x - y)g(y)\eu^{-\iu\xi x}$ \omterm{def:b3:lebesgue:l1}{terintegralkan}
mutlak (\cref{thm:b3:product:convolution}):
\[
\widehat{f*g}(\xi) = \iint f(x - y)g(y)\eu^{-\iu\xi(x - y)}
\eu^{-\iu\xi y}\dd x\,\dd y = \hat f(\xi)\,\hat g(\xi).
\]
\end{proof}

\begin{example}\label{ex:b3:fouriertransform:gaussian}
Fungsi Gauss: untuk $a > 0$,\index{fungsi Gauss, transformasi Fourier}
\[
\widehat{\eu^{-ax^2}}(\xi) =
\sqrt{\frac\pi a}\;\eu^{-\xi^2/4a} :
\]
menurut \cref{exo:b3:lebesgue:7} (lewat kiat persamaan diferensial $F' = -\frac\xi{2}F$
yang diskalakan ulang), atau lewat (3): sebab $g = \widehat{\eu^{-ax^2}}$
memenuhi $g'(\xi) = -\frac{\xi}{2a}g(\xi)$ (lewat pengintegralan
parsial) dengan $g(0) = \sqrt{\pi/a}$. Jadi fungsi Gauss merupakan titik tetap
transformasinya sampai penskalaan --- yakni alasan mendalam mereka menguasai
teorema limit pusat.
\end{example}

\begin{theorem}[Riemann--Lebesgue]
\label{thm:b3:fouriertransform:riemannlebesgue}
Untuk $f \in L^1(\R)$: $\hat f(\xi) \to 0$ saat $\abs\xi \to
\infty$. Jadi $\widehat{\phantom f} \colon L^1 \to \mathcal
C_0(\R)$ (yakni fungsi \omterm{def:b3:topology:continuity}{kontinu} yang lenyap di
tak hingga).\index{lema Riemann--Lebesgue}
\end{theorem}

\begin{proof}
Untuk indikator sebuah selang, $\widehat{\mathbf 1_{\intcc
ab}}(\xi) = \frac{\eu^{-\iu a\xi} - \eu^{-\iu b\xi}}{\iu\xi}
\to 0$; sehingga demikian pula untuk fungsi tangga. Lalu fungsi tangganya padat di
$L^1$ (\cref{thm:b3:lp:density}(1) ditambah penghampiran
himpunan berukuran berhingga oleh gabungan berhingga selang,
\cref{exo:b3:measure:7}), dan transformasinya bersifat
\omterm{def:b3:topology:continuity}{kontinu} $\norm\cdot_\infty$-$\norm\cdot_1$: sehingga untuk $\norm{f -
s}_1 < \varepsilon$, $\limsup_{\abs\xi\to\infty}\abs{\hat
f(\xi)} \leq \varepsilon$.
\end{proof}

\section{Pembalikan dan keinjektifan}

\begin{lemma}[Rumus perkalian]
\label{lem:b3:fouriertransform:multiplication}
Untuk $f, g \in L^1(\R)$: $\displaystyle\int \hat f\,g =
\int f\,\hat g$.
\end{lemma}

\begin{proof}
Kedua ruasnya sama dengan $\iint f(x)g(\xi)\eu^{-\iu x\xi}\dd
x\,\dd\xi$ (lewat Tonelli--Fubini: sebab integral rangkap atas
nilai mutlaknya adalah $\norm f_1\norm g_1$).
\end{proof}

\begin{theorem}[Pembalikan]\label{thm:b3:fouriertransform:inversion}
Misalkan $f \in L^1(\R)$.\index{teorema pembalikan Fourier}
\begin{enumerate}
  \item (Keterjumlahan Gauss) Untuk setiap $x$,
        \[
        (f * g_\varepsilon)(x) =
        \frac1{2\pi}\int_\R \hat f(\xi)\,
        \eu^{-\varepsilon\xi^2}\,\eu^{\iu x\xi}\,\dd\xi,
        \qquad\text{dengan } g_\varepsilon(y) =
        \frac{1}{2\sqrt{\pi\varepsilon}}\,
        \eu^{-y^2/4\varepsilon},
        \]
        dan $f * g_\varepsilon \to f$ di $L^1$ saat $\varepsilon
        \to 0$.
  \item Jika lebih lanjut $\hat f \in L^1$, maka untuk hampir setiap
        $x$
        \[
        f(x) = \frac{1}{2\pi}\int_\R \hat
        f(\xi)\,\eu^{\iu x\xi}\,\dd\xi ,
        \]
        dan $f$ mempunyai wakil yang \omterm{def:b3:topology:continuity}{kontinu}.
  \item (Keinjektifan) Jika $\hat f = 0$ maka $f = 0$ hampir di mana-mana.
\end{enumerate}
\end{theorem}

\begin{proof}
(1) Tetapkan $x$ lalu terapkan
\cref{lem:b3:fouriertransform:multiplication} pada $f$ dan
$g(\xi) = \frac1{2\pi}\eu^{-\varepsilon\xi^2}\eu^{\iu x\xi}$:
maka menurut \cref{ex:b3:fouriertransform:gaussian} (dengan aturan
modulasinya),
\[
\hat g(y) = \frac1{2\pi}\sqrt{\frac\pi\varepsilon}\,
\eu^{-(y - x)^2/4\varepsilon} = g_\varepsilon(x - y),
\]
sehingga $\frac1{2\pi}\int\hat f(\xi)\eu^{-\varepsilon\xi^2}
\eu^{\iu x\xi}\dd\xi = \int f(y)g_\varepsilon(x - y)\dd y =
(f*g_\varepsilon)(x)$. Sedangkan $g_\varepsilon$ merupakan
identitas hampiran: sebab $g_\varepsilon \geq 0$,
$\int g_\varepsilon = 1$ (lewat integral Gaussnya), dan ia memusat di
$0$; sehingga bukti \cref{thm:b3:lp:regularization}(2) berlaku
kata demi kata (sebab hanya $\int g_\varepsilon = 1$ dan pemusatannya yang
dipakai: dan untuk ekornya, $\int_{\abs y > \delta}g_\varepsilon \to
0$): jadi $\norm{f * g_\varepsilon - f}_1 \to 0$.

(2) Jika $\hat f \in L^1$: maka ruas kanan (1) konvergen, lewat
kekonvergenan terdominasi (dengan pendominasi $\abs{\hat f}$), ke $\frac1{2\pi}\int\hat
f(\xi)\eu^{\iu x\xi}\dd\xi$ untuk \emph{setiap} $x$, dan fungsi
limitnya \omterm{def:b3:topology:continuity}{kontinu} (lewat kekonvergenan terdominasi lagi). Sedangkan di sisi lain $f
* g_\varepsilon \to f$ di $L^1$, sehingga sepanjang sebuah subbarisan hampir di mana-mana
(\cref{thm:b3:lp:complete}): jadi kedua limitnya bersesuaian hampir di mana-mana.

(3) Karena $\hat f = 0$ membuat ruas kanan (1) lenyap: $f *
g_\varepsilon = 0$ untuk setiap $\varepsilon$, dan $f *
g_\varepsilon \to f$ di $L^1$: sehingga $f = 0$ hampir di mana-mana.
\end{proof}

\section{Kelas Schwartz}

\begin{definition}\label{def:b3:fouriertransform:schwartz}
\emph{Kelas Schwartz}\index{kelas Schwartz} $\mathcal
S(\R)$ terdiri atas fungsi $\mathcal C^\infty$ yang $f$-nya memenuhi
$\sup_x\abs{x^m f^{(n)}(x)} < \infty$ untuk setiap $m, n \geq 0$
(sehingga semua turunannya meluruh lebih cepat daripada pangkat mana pun). Contohnya:
$\eu^{-ax^2}$ dan $\mathcal C_c^\infty$. Jelaslah $\mathcal S
\subseteq L^p$ untuk setiap $p$ (lewat batas $C(1 + x^2)^{-1}$), dan
$\mathcal S$ stabil terhadap penurunan, perkalian dengan
polinomial, dan hasil kali.
\end{definition}

\begin{theorem}\label{thm:b3:fouriertransform:schwartz}
Transformasi Fourier memetakan $\mathcal S(\R)$ secara bijektif ke
dirinya sendiri, dengan balikan $\check g(x) = \frac1{2\pi}\int
g(\xi)\eu^{\iu x\xi}\dd\xi$.
\end{theorem}

\begin{proof}
Misalkan $f \in \mathcal S$. Dengan mengiterasikan
\cref{prop:b3:fouriertransform:basic}(3), $(\hat f)^{(n)} =
\widehat{(-\iu x)^nf}$ (sebab setiap $x^kf \in L^1$); lalu mengiterasikan (4)
dengan $h = (-\iu x)^nf \in \mathcal S$ (yang semua
turunannya \omterm{def:b3:lebesgue:l1}{terintegralkan}), $(\iu\xi)^m\hat h =
\widehat{h^{(m)}}$. Digabungkan,
\[
\abs{\xi^m\,(\hat f)^{(n)}(\xi)}
= \bigl|\widehat{\,h^{(m)}}(\xi)\bigr|
\leq \bigl\|\bigl((-\iu x)^nf\bigr)^{(m)}\bigr\|_1 < \infty
\]
secara seragam terhadap $\xi$: jadi $\hat f \in \mathcal S$. Lalu karena $\hat f \in L^1$, pembalikannya
(\cref{thm:b3:fouriertransform:inversion}(2)) berlaku
di mana-mana (sebab kedua ruasnya \omterm{def:b3:topology:continuity}{kontinu}): sehingga $\check{\hat f} = f$, dan
setangkup dengan itu $\widehat{\check g} = g$ (sebab transformasi centangnya
adalah $g \mapsto \frac1{2\pi}\hat g(-\cdot)$, yang kembali mengawetkan
$\mathcal S$): jadi bijeksi.
\end{proof}

\section{Plancherel dan \texorpdfstring{$L^2$}{L2}}

\begin{theorem}[Plancherel]\label{thm:b3:fouriertransform:plancherel}
Untuk $f \in L^1 \cap L^2(\R)$:\index{teorema Plancherel}
\[
\norm{\hat f}_2^2 = 2\pi\,\norm f_2^2 .
\]
Karena itu $\widehat{\phantom f}$ diperluas secara tunggal menjadi sebuah
pemetaan linear \omterm{def:b3:topology:continuity}{kontinu} $\mathcal F \colon L^2(\R) \to L^2(\R)$
dengan $\norm{\mathcal Ff}_2 = \sqrt{2\pi}\norm f_2$; dan $\mathcal
F$ bersifat bijektif, dengan $\mathcal F^{-1} =
\frac1{2\pi}\,\mathcal F\circ\sigma$ dengan $\sigma f =
f(-\cdot)$, serta ia mengawetkan \omterm{def:b3:hilbert:inner}{hasil kali dalam} sampai faktor
$2\pi$.
\end{theorem}

\begin{proof}
Misalkan $f \in L^1\cap L^2$ dan $h = f * \tilde f$ dengan
$\tilde f(x) = \overline{f(-x)}$. Maka $h \in L^1$
(\cref{thm:b3:product:convolution}), $h$ \omterm{def:b3:topology:continuity}{kontinu} dan
terbatas (\cref{exo:b3:lp:6}: sebab $f, \tilde f \in L^2$), $h(0) =
\int f\bar f = \norm f_2^2$, dan $\hat h = \hat
f\,\widehat{\tilde f} = \hat f\,\overline{\hat f} =
\abs{\hat f}^2 \geq 0$ (hitunglah $\widehat{\tilde f} =
\overline{\hat f}$). Lalu terapkan
\cref{thm:b3:fouriertransform:inversion}(1) pada $h$ di $x = 0$:
\[
(h * g_\varepsilon)(0) = \frac1{2\pi}\int \hat
h(\xi)\,\eu^{-\varepsilon\xi^2}\dd\xi .
\]
Saat $\varepsilon \to 0$: ruas kirinya menuju $h(0)$ (sebab $h$
\omterm{def:b3:topology:continuity}{kontinu} terbatas: $(h*g_\varepsilon)(0) - h(0) = \int(h(-y)
- h(0))g_\varepsilon(y)\dd y \to 0$ lewat pemisahan $y$ yang kecil
dan besar); sedangkan ruas kanannya naik menuju $\frac1{2\pi}\int\hat h$ lewat
kekonvergenan monoton (sebab $\hat h \geq 0$). Karena itu $\frac1{2\pi}\int\abs{\hat f}^2 =
\norm f_2^2$, entah berhingga entah tidak \emph{a priori} --- dan ternyata berhingga,
yang membuktikan keanggotaannya sekaligus identitasnya.

Untuk perluasannya: $L^1\cap L^2 \supseteq \mathcal C_c$ bersifat padat di
$L^2$ (\cref{thm:b3:lp:density}); transformasinya bersifat
isometrik sampai $\sqrt{2\pi}$ di sana, sehingga diperluas secara tunggal menjadi
isometri sampai konstanta $\mathcal F$ pada $L^2$
(\cref{thm:b3:complete:extension}). Lalu pembalikan pada $\mathcal S$
(\cref{thm:b3:fouriertransform:schwartz}) berpindah lewat
\omterm{ex:b3:lebesgue:gamma}{kepadatan} yang sama (sebab kedua ruasnya kontinu-$L^2$): $\mathcal
F\bigl(\frac1{2\pi}\mathcal F(\sigma f)\bigr) = f$ pada
$\mathcal S$, jadi pada $L^2$: sehingga bijektif. Untuk \omterm{def:b3:hilbert:inner}{hasil kali dalamnya}:
lewat polarisasi dari identitas normanya.
\end{proof}

\begin{theorem}[Penjumlahan Poisson]
\label{thm:b3:fouriertransform:poisson}
Misalkan $f \in \mathcal S(\R)$ (sebab $f$ \omterm{def:b3:topology:continuity}{kontinu} dengan $\abs{f} +
\abs{\hat f} \leq C(1 + \abs\cdot)^{-2}$
sudah cukup). Maka\index{rumus penjumlahan Poisson}
\[
\sum_{n\in\Z} f(n) \;=\; \sum_{k\in\Z}\hat f(2\pi k) .
\]
\end{theorem}

\begin{proof}
Misalkan $F(x) = \sum_{n\in\Z}f(x + n)$: deretnya konvergen
secara normal pada \omterm{def:b3:topology:compact}{kompak} (lewat peluruhan $f$), sehingga $F$ \omterm{def:b3:topology:continuity}{kontinu}, dan
ia berperiode $1$. Koefisien Fouriernya (dengan periode $1$:
$c_k(F) = \int_0^1F(t)\eu^{-2\iu\pi kt}\dd t$):
\[
c_k(F) = \sum_n\int_0^1 f(t + n)\,\eu^{-2\iu\pi kt}\dd t
= \int_\R f(t)\,\eu^{-2\iu\pi kt}\dd t = \hat f(2\pi k)
\]
(sebab kekonvergenan normalnya membenarkan penukarannya; dan fasanya
berperiode $1$). Deret $\sum_k\abs{c_k(F)}$ konvergen
(lewat peluruhan $\hat f$), sehingga deret Fourier $F$ konvergen
secara normal; jumlahnya berupa fungsi \omterm{def:b3:topology:continuity}{kontinu} yang berkoefisien
Fourier sama dengan $F$, jadi sama dengan $F$ (lewat keinjektifan pada
lingkarannya: sebab selisihnya berkoefisien nol, dan
\cref{thm:b3:hilbert:fourier} memberikan nol di $L^2$, sehingga
di mana-mana lewat \omterm{def:b3:topology:continuity}{kekontinuannya}). Lalu nilaikan di $x = 0$.
\end{proof}

\begin{example}[Identitas theta]
\label{ex:b3:fouriertransform:theta}
Dengan menerapkan Poisson pada $f(x) = \eu^{-\pi tx^2}$ (dengan $t > 0$), yang
transformasinya $\hat f(\xi) = t^{-1/2}\eu^{-\xi^2/4\pi t}$
(\cref{ex:b3:fouriertransform:gaussian} dengan $a = \pi t$):
\[
\sum_{n\in\Z}\eu^{-\pi n^2t}
= \frac1{\sqrt t}\sum_{k\in\Z}\eu^{-\pi k^2/t} :
\]
yakni persamaan fungsional fungsi theta Jacobi, kunci bagi
persamaan fungsional $\zeta$ Riemann --- sekaligus pemercepat
numerik yang spektakuler: sebab untuk $t$ yang kecil, ruas kirinya konvergen
lambat, sedangkan ruas kanannya sangat cepat.
\end{example}

\begin{method}\label{met:b3:fouriertransform:method}
Rentang kerjanya: $L^1$ --- transformasinya terdefinisi titik demi titik,
sedangkan pembalikannya menuntut $\hat f \in L^1$; $\mathcal S$ --- segalanya
sah, jadi buktikanlah di sini dahulu; $L^2$ --- transformasinya terdefinisi lewat
\omterm{ex:b3:lebesgue:gamma}{kepadatan} (bukan lewat integralnya!), dengan kesetangkupan \omterm{prop:b3:galois:perfect}{sempurna} dan
pembukuan \omterm{thm:b3:hilbert:parseval}{Parseval}. Untuk menghitung sebuah transformasi: susutkanlah ke tabelnya
(indikator, eksponensial, fungsi Gauss) lewat aturan
\cref{prop:b3:fouriertransform:basic}; untuk membuktikan sebuah identitas:
tegakkanlah pada $\mathcal S$ (atau $\mathcal C_c^\infty$) lalu
perluas lewat \omterm{ex:b3:lebesgue:gamma}{kepadatan} dan \omterm{def:b3:topology:continuity}{kekontinuannya}
(\cref{met:b3:lp:toolkit}); dan untuk menyelesaikan persamaan diferensial parsial atau biasa yang
berkoefisien konstan: transformasikan, bagilah, lalu balikkan.
\end{method}

\begin{omfigure}
\begin{tikzpicture}[xscale=1.05, yscale=2.2]
  \draw[->] (-4.3,0) -- (4.3,0) node[right,font=\small] {$x$};
  \draw[->] (0,0) -- (0,1.25);
  \draw[omThm, thick, smooth]
    plot[domain=-4:4, samples=81]
    (\x, {exp(-\x*\x/0.4)/1.1209});
  \draw[omProp, thick, smooth]
    plot[domain=-4:4, samples=81]
    (\x, {exp(-\x*\x/1.6)/2.2418});
  \draw[omDef, thick, smooth]
    plot[domain=-4:4, samples=81]
    (\x, {exp(-\x*\x/6.4)/4.4836});
  \node[omThm, font=\small] at (1.35,0.75) {$t = 0.1$};
  \node[omProp, font=\small] at (2.4,0.36) {$t = 0.4$};
  \node[omDef, font=\small] at (3.5,0.17) {$t = 1.6$};
\end{tikzpicture}

{\small \omterm{pb:b3:fouriertransform:1}{Kernel panas} $g_t(x) =
\frac1{2\sqrt{\pi t}}\,\eu^{-x^2/4t}$ pada tiga waktu: dengan massa
total $1$ selamanya, tinggi $\sim t^{-1/2}$, dan lebar $\sim
\sqrt t$. Mengonvolusikan data awal dengan fungsi Gauss yang
menyebar ini adalah seluruh isi soal akhir pekannya; sedangkan dalam
frekuensinya, gambar yang sama terbaca
$\hat g_t(\xi) = \eu^{-t\xi^2}$ --- yakni frekuensi tinggi mati
lebih dahulu, dan ketaksetangkupan itulah anak panah waktu.}
\end{omfigure}

\section{Latihan}

\begin{exercise}[$\star$]\label{exo:b3:fouriertransform:1}
Hitunglah transformasi Fourier dari: $\mathbf 1_{\intcc{-a}a}$;
$\eu^{-a\abs x}$ (dengan $a > 0$); fungsi tenda $\max(0, 1 -
\abs x)$; dan $\frac1{x^2 + a^2}$ \emph{(pakailah pembalikan pada
yang kedua)}. Lalu catatlah tabel yang muncul.
\end{exercise}

\begin{exercise}[$\star$]\label{exo:b3:fouriertransform:2}
Misalkan $f \in L^1$. Nyatakan dalam $\hat f$ transformasi
dari: $f(x - a)$, $f(x)\cos(bx)$, $f(ax + b)$,
$\overline{f(-x)}$, dan $(f * f)(x)$. Lalu periksalah setiap aturannya pada
fungsi Gauss.
\end{exercise}

\begin{exercise}[$\star\star$]\label{exo:b3:fouriertransform:3}
(a) Tunjukkan bahwa $\mathbf 1_{\intcc{-1}1} * \mathbf
1_{\intcc{-1}1}$ bertransformasi $\bigl(\frac{2\sin\xi}\xi
\bigr)^2$, lalu turunkan
$\int_\R\bigl(\frac{\sin\xi}\xi\bigr)^2\dd\xi = \pi$ lewat
Plancherel --- atau lewat pembalikan di $0$. Bandingkan dengan
\cref{pb:b3:lebesgue:1}.
(b) Hitunglah $\int_\R\frac{\dd x}{(x^2+1)^2}$ lewat Plancherel
yang diterapkan pada $\eu^{-\abs x}$.
\end{exercise}

\begin{exercise}[$\star\star$]\label{exo:b3:fouriertransform:4}
(Aljabar \omterm{pb:b3:fouriertransform:1}{kernel panas}) Dengan $g_t(x) =
\frac1{2\sqrt{\pi t}}\eu^{-x^2/4t}$:
(a) periksalah $\hat g_t(\xi) = \eu^{-t\xi^2}$;
(b) turunkan hukum semigrupnya $g_t * g_s = g_{t+s}$ tanpa
perhitungan integral apa pun;
(c) tunjukkan $\norm{g_t}_1 = 1$ dan $\norm{g_t}_2^2 =
(8\pi t)^{-1/2}$.
\end{exercise}

\begin{exercise}[$\star\star$]\label{exo:b3:fouriertransform:5}
Tunjukkan bahwa jika $f \in L^1$ genap dan real, maka $\hat f$ genap
dan real; sedangkan jika $f$ gasal dan real, maka $\hat f$ gasal dan murni
imajiner. Apakah yang dihitung $\hat f(0)$? Lalu turunkan bahwa $f \geq
0$ memaksa $\norm{\hat f}_\infty = \hat f(0) = \int f$, lalu
tafsirkanlah untuk \omterm{ex:b3:lebesgue:gamma}{kepadatan} peluang (\cref{ch:b3:clt}: sebab sebuah
fungsi karakteristik bermodulus $\leq 1$, yang tercapai di
$0$).
\end{exercise}

\begin{exercise}[$\star\star\star$]\label{exo:b3:fouriertransform:6}
(Ketaksurjektifan) Tunjukkan bahwa $\widehat{\phantom f}\colon L^1
\to \mathcal C_0$ bersifat injektif dan \omterm{def:b3:topology:continuity}{kontinu}, tetapi \emph{tak}
surjektif, dalam tiga langkah.
(i) Keinjektifannya (\cref{thm:b3:fouriertransform:inversion}) dan
\omterm{def:b3:topology:continuity}{kekontinuannya} (sebab $\norm{\hat f}_\infty \leq \norm f_1$), dan
$\mathcal C_0$ merupakan ruang Banach (yang tertutup dalam
$\norm\cdot_\infty$).
(ii) Seandainya pemetaannya surjektif, ia akan bijektif, dan
teorema pemetaan terbuka (\cref{thm:b3:banach:openmapping})
akan memberikan sebuah konstanta $C$ dengan $\norm f_1 \leq C\norm{\hat
f}_\infty$ untuk setiap $f \in L^1$.
(iii) Pertentangkanlah itu dengan $f_n(x) = \frac{\sin
x}{x}\cdot\frac{\sin(x/n)}{x/n}$: sebab transformasinya berupa (sampai
konstantanya) trapesium bertipe konvolusi $\mathbf 1_{\intcc{-1}1} *
\mathbf 1_{\intcc{-1/n}{1/n}}$ --- lalu tunjukkan
$\norm{\hat f_n}_\infty \leq \pi$ secara seragam, sedangkan
$\norm{f_n}_1 \geq c\ln n$ dengan mencacah lengkungan
$\frac{\abs{\sin x}}x$ pada $[1, n]$ (tempat faktor keduanya
terbatas dari bawah), seperti pada \cref{thm:b3:banach:fourierdiverge}.
\end{exercise}

\begin{exercise}[$\star\star$]\label{exo:b3:fouriertransform:7}
(Kamus kemulusan $\leftrightarrow$ peluruhan) Buktikan: bahwa $f \in
L^1$ dengan $\hat f(\xi) = O(\abs\xi^{-k-1-\delta})$ untuk suatu
$\delta > 0$ mengakibatkan $f$ mempunyai wakil $\mathcal C^k$.
Sebaliknya $f \in \mathcal C^k_c$ mengakibatkan $\hat f(\xi) =
O(\abs\xi^{-k})$. Lalu peragakan kedua arahnya pada fungsi
tenda.
\end{exercise}

\begin{exercise}[$\star\star\star$]\label{exo:b3:fouriertransform:8}
(Ketaksamaan Heisenberg) Untuk $f \in \mathcal S(\R)$ yang real dengan
$\norm f_2 = 1$, buktikan bahwa
\[
\Bigl(\int x^2f(x)^2\dd x\Bigr)\cdot
\Bigl(\frac1{2\pi}\int \xi^2\abs{\hat f(\xi)}^2\dd\xi\Bigr)
\;\geq\; \frac14 ,
\]
dengan kesamaannya bagi fungsi Gauss. \emph{(Tulislah $1 = \int f^2 =
-\int x\,(f^2)'$ lewat pengintegralan parsial, batasi lewat Cauchy--Schwarz, lalu
ubahlah $\norm{f'}_2$ lewat Plancherel.)} Tafsirannya: bahwa sebuah
sinyal dan spektrumnya tak dapat sama-sama terpusat.
\end{exercise}

\begin{exercise}[$\star\star$]\label{exo:b3:fouriertransform:9}
Benarkanlah \cref{ex:b3:fouriertransform:theta} secara terperinci
(yakni hipotesis Poisson bagi fungsi Gaussnya), lalu pakailah identitasnya
untuk menilaikan $\sum_{n\in\Z}\eu^{-\pi n^2}$ sampai enam desimal dengan
tiga suku. Berapa suku deret \emph{pendefinisinya}
yang diperlukan bagi ketelitian yang sama pada $t = 10^{-2}$, dibandingkan dengan
deret yang tertransformasi?
\end{exercise}

\begin{exercise}[$\star\star$]\label{exo:b3:fouriertransform:10}
(Fungsi terbatas pita) Misalkan $f \in L^2(\R)$ dengan $\mathcal Ff$
berpendukung di $\intcc{-\pi}\pi$. Tunjukkan bahwa $f$ mempunyai sebuah
wakil yang setiap nilainya dapat dipulihkan
dari cuplikannya: buktikanlah \emph{interpolasi Shannon} di
bilangan bulatnya,
\[
f(x) = \sum_{n\in\Z} f(n)\,
\frac{\sin\bigl(\pi(x - n)\bigr)}{\pi(x - n)}
\quad\text{di } L^2,
\]
dengan menguraikan $\mathcal Ff$ pada basis Fourier
$L^2(\intcc{-\pi}\pi)$ (\cref{thm:b3:hilbert:fourier}) lalu
mentransformasikannya kembali suku demi suku.
\end{exercise}

\begin{exercise}[$\star\star$]\label{exo:b3:fouriertransform:11}
(Transformasinya sebagai operator berorde empat) Pada $\mathcal
S(\R)$, misalkan $\mathcal F f = \hat f$.
(a) Dengan memakai rumus pembalikannya, tunjukkan $(\mathcal F^2f)(x) =
2\pi\,f(-x)$, lalu turunkan $\mathcal F^4 = (2\pi)^2\,
\mathrm{id}$.
(b) Turunkan bahwa setiap nilai eigen $\mathcal F$ pada
$\mathcal S$ termasuk $\{\pm\sqrt{2\pi},
\pm\iu\sqrt{2\pi}\}$, lalu tampilkan sebuah fungsi eigen untuk
$+\sqrt{2\pi}$ \emph{(fungsi mana pada bab ini yang
sebanding dengan transformasinya sendiri?)}.
(c) Tunjukkan bahwa fungsi genap memenuhi $\mathcal F^2f = 2\pi
f$ dan yang gasal $\mathcal F^2f = -2\pi f$; lalu hasilkanlah sebuah
fungsi eigen untuk nilai eigen $-\iu\sqrt{2\pi}$ dari
$x\eu^{-x^2/2}$ dengan menghitung transformasinya (turunkanlah
transformasi fungsi Gaussnya).
\end{exercise}

\begin{exercise}[$\star\star$]\label{exo:b3:fouriertransform:12}
(Autokorelasi dan lema Wiener) Untuk $f \in L^2(\R)$
definisikan $\tilde f(x) = \overline{f(-x)}$ dan
\emph{autokorelasinya} $A_f = f * \tilde f$.
(a) Tunjukkan bahwa $A_f$ merupakan fungsi \omterm{def:b3:topology:continuity}{kontinu} terbatas dengan
$A_f(0) = \norm f_2^2 \geq \abs{A_f(x)}$ untuk setiap $x$
(lewat \cref{exo:b3:lp:6} dan Cauchy--Schwarz).
(b) Tunjukkan, mula-mula untuk $f \in L^1\cap L^2$, bahwa
$\widehat{A_f} = \abs{\hat f\,}^2 \geq 0$: sehingga
autokorelasinya bertransformasi tak negatif --- jadi spektrum
autokorelasi merupakan spektrum daya.
(c) Turunkan identitas $\int_\R\abs{\hat
f(\xi)}^2\eu^{\iu x\xi}\,\dd\xi = 2\pi A_f(x)$ (lewat pembalikan;
benarkanlah keterterapannya bila $\hat f \in L^2$ mempunyai
$\abs{\hat f}^2 \in L^1$), lalu nilaikanlah untuk $f = \mathbf
1_{\intcc{-1/2}{1/2}}$ di $x = 0$: sehingga pulihlah
$\int_\R\bigl(\frac{\sin u}u\bigr)^2\dd u = \pi$.
\end{exercise}

\section{Soal: persamaan panas pada garis}

\begin{problem}[{Soal akhir pekan --- $\partial_tu =
\partial^2_{xx}u$, yang diselesaikan sampai tuntas}]
\label{pb:b3:fouriertransform:1}
Panas menyebar; dan persamaan $\partial_tu = \partial_{xx}^2u$
mengatakan bahwa \omterm{ex:b3:lebesgue:gamma}{kepadatannya} berdifusi dengan laju yang diberikan kelengkungan
lokal profil suhunya. Kita menyelesaikan soal Cauchynya
pada $\R$ --- yakni diberikan $f$, carilah $u(t, x)$ untuk $t > 0$ dengan
$u(0, \cdot) = f$ --- lalu membuktikan sifat solusinya yang
luar biasa, dan melihat mengapa waktu tak dapat dibalik. Sepanjang soal ini,
$g_t(x) = \frac1{2\sqrt{\pi t}}\eu^{-x^2/4t}$ adalah
\emph{kernel panas}\index{kernel panas} dan $u(t, \cdot) = g_t *
f$.

\medskip
\noindent\textbf{Bagian I --- Menurunkan kernelnya.} Bekerjalah
secara formal dahulu: andaikan $u(t, \cdot) \in \mathcal S$ menyelesaikan
persamaannya, lalu misalkan $\hat u(t, \xi)$ transformasinya terhadap
$x$.
\begin{enumerate}
  \item Tunjukkan (secara formal) $\partial_t\hat u = -\xi^2\hat u$,
        sehingga $\hat u(t,\xi) = \eu^{-t\xi^2}\hat f(\xi)$, lalu
        kenalilah $u(t) = g_t * f$
        (\cref{exo:b3:fouriertransform:4}). Ini memotivasi
        \emph{definisi} $u$; dan segalanya kini dibuktikan
        secara langsung, untuk $f \in \mathcal C_b(\R)$ (yakni \omterm{def:b3:topology:continuity}{kontinu}
        terbatas) atau $f \in L^p$.
\end{enumerate}

\medskip
\noindent\textbf{Bagian II --- Pemeriksaannya.}
\begin{enumerate}[resume]
  \item Tunjukkan bahwa untuk $t > 0$, $u(t, x) = \int
        g_t(x-y)f(y)\dd y$ terdefinisi dengan baik untuk $f \in \mathcal
        C_b$, dan bahwa $u$ bersifat $\mathcal C^\infty$ terhadap $(t, x)$
        pada $\intoo0\infty\times\R$ \emph{(turunkanlah di bawah
        integralnya; lalu dominasi turunan Gaussnya secara seragam lokal terhadap $(t,x)$)}.
  \item Periksalah $\partial_tg_t = \partial^2_{xx}g_t$ lewat perhitungan
        langsung, lalu turunkan $\partial_tu =
        \partial^2_{xx}u$ untuk $t > 0$.
  \item (Syarat awalnya) Tunjukkan bahwa untuk $f \in \mathcal
        C_b$, $u(t, x) \to f(x)$ saat $t \to 0^+$, secara seragam pada
        \omterm{def:b3:topology:compact}{kompak} \emph{(lewat identitas hampirannya: pisahkan $\abs y
        \leq \delta$ dan $\abs y > \delta$)}; sedangkan untuk $f \in L^p$
        (dengan $p < \infty$), tunjukkan $\norm{u(t) - f}_p \to 0$.
  \item (Pemulusan seketika) Simpulkan: bahkan untuk $f$ yang sekadar
        \omterm{def:b3:topology:continuity}{kontinu} terbatas, solusinya bersifat $\mathcal C^\infty$
        untuk setiap $t > 0$ --- sehingga panas seketika menghapus
        kekasarannya. Lalu hitunglah $u(t, \cdot)$ secara eksplisit untuk $f =
        \mathbf 1_{\intoo0\infty}$ (yakni sebuah fungsi galat) lalu
        sketsakan profilnya untuk tiga nilai $t$.
\end{enumerate}

\medskip
\noindent\textbf{Bagian III --- Sifat strukturalnya.}
\begin{enumerate}[resume]
  \item (Kepositifan dan pembandingan) Jika $f \geq 0$ maka $u > 0$
        untuk setiap $t > 0$ (secara sejati, kecuali $f = 0$ hampir di mana-mana); dan jika
        $f_1 \leq f_2$ maka $u_1 \leq u_2$. Sebuah titik dingin menghangat
        seketika: berkomentarlah.
  \item (Kekekalan) Untuk $f \in L^1$: $\int u(t, x)\dd x =
        \int f$ untuk setiap $t$ \emph{(lewat Tonelli)} --- sehingga panas totalnya
        kekal.
  \item (Pelesapan) Untuk $f \in L^1\cap L^2$, tunjukkan lewat
        Plancherel bahwa $t \mapsto \norm{u(t)}_2$ bersifat
        tak naik, secara sejati kecuali $f = 0$, lalu hitunglah
        limitnya saat $t \to \infty$. Lalu tunjukkan lebih lanjut
        $\norm{u(t)}_\infty \leq \frac{\norm
        f_1}{2\sqrt{\pi t}} \to 0$: sehingga panas menyebar lalu mati.
  \item (Ketunggalan, kelas $L^2$) Misalkan $u$ sebuah solusi dengan
        $u(t) \in L^2$ untuk setiap $t$, $u \in \mathcal
        C^1(\intoo0\infty, L^2)$ dalam arti yang wajar, dan
        $u(t) \to f$ di $L^2$ saat $t\to0$; lalu dengan menerima bahwa
        transformasinya mengubahnya menjadi $\partial_t\hat u =
        -\xi^2\hat u$ secara titik demi titik hampir di mana-mana pada $\xi$ untuk hampir setiap $t$
        \emph{(yang dibenarkan lewat pengujian terhadap $\mathcal
        C_c^\infty$ pada $\xi$ --- garis besarkanlah ini)}, tunjukkan
        $\hat u(t,\xi) = \eu^{-t\xi^2}\hat f(\xi)$, sehingga
        ketunggalannya pada kelas ini.
\end{enumerate}

\medskip
\noindent\textbf{Bagian IV --- Anak panah waktu.}
\begin{enumerate}[resume]
  \item Tunjukkan bahwa soal \emph{mundurnya} tak berkondisi baik:
        yakni agar solusinya ada pada waktu $-s$ (dengan $s > 0$) dengan
        data $f$ pada waktu $0$ --- yakni agar $f = g_s * h$
        mempunyai solusi $h \in L^2$ --- diperlukan bahwa
        $\eu^{s\xi^2}\hat f(\xi) \in L^2$: yakni syarat peluruhan yang
        ekstrem pada $\hat f$. Tampilkan sebuah $f \in L^2$ yang mulus dan
        eksplisit yang tak bersolusi mundur pada selang
        waktu mana pun: ambillah fungsi yang $\hat f(\xi) =
        \eu^{-\abs\xi}$-nya --- kenalilah $f$
        (\cref{exo:b3:fouriertransform:1}) lalu tunjukkan
        $\eu^{s\xi^2}\eu^{-\abs\xi} \notin L^2$ untuk setiap $s >
        0$.
  \item (Pemulusan lawan keterangan) Terangkan dalam sebuah paragraf
        pendek, dengan memakai pertanyaan 5, 9, dan 10, mengapa semigrup
        panasnya $(f \mapsto g_t * f)_{t\geq0}$ bersifat injektif
        tetapi tak surjektif pada $L^2$, dan mengapa ini menyatakan
        ketakterbalikan difusinya.
\end{enumerate}

\medskip
\noindent\textbf{Bagian V --- Teorema pencuplikan Shannon.} Sebuah
fungsi $f \in L^2(\R)$ disebut \emph{terbatas pita} pada $\Omega$
jika $\hat f = 0$ hampir di mana-mana di luar $\intcc{-\Omega}\Omega$; lalu tulislah
$PW_\Omega$ (yakni ruang Paley--Wiener) untuk fungsi semacam itu.
\begin{enumerate}[resume]
  \item Tunjukkan bahwa setiap $f \in PW_\Omega$ sama hampir di mana-mana dengan
        fungsi $\mathcal C^\infty$
        $\frac1{2\pi}\int_{-\Omega}^{\Omega}\hat
        f(\xi)\eu^{\iu x\xi}\,\dd\xi$ (benarkanlah kemulusannya dan
        pengenalannya), dengan semua turunannya terbatas:
        sehingga pembatasan pita merupakan bentuk keteraturan yang ekstrem.
        Mulai sekarang $f$ menandai wakil itu.
  \item Uraikan $\hat f \in L^2(\intcc{-\Omega}\Omega)$ pada
        basis Fourier selang itu lalu kenalilah
        koefisiennya sebagai \emph{cuplikan} $f$:
        \[
        \hat f(\xi) = \frac\pi\Omega\sum_{n\in\Z}
        f\Bigl(\frac{n\pi}\Omega\Bigr)\,
        \eu^{-\iu n\pi\xi/\Omega}
        \quad\text{di } L^2(\intcc{-\Omega}\Omega) .
        \]
  \item Turunkan \emph{teorema pencuplikannya}: bahwa untuk $f \in
        PW_\Omega$,
        \[
        f(x) = \sum_{n\in\Z}f\Bigl(\frac{n\pi}
        \Omega\Bigr)\,\operatorname{sinc}(\Omega x - n\pi),
        \qquad \operatorname{sinc}t = \frac{\sin t}t,
        \]
        dengan kekonvergenan di $L^2(\R)$ dan secara seragam pada $\R$
        \emph{(sisipkan deret pertanyaan 13 ke dalam
        rumus pembalikannya lalu hitunglah integral
        dasarnya)}: sehingga sinyal terbatas pita sepenuhnya
        ditentukan oleh nilainya pada kisi berlangkah
        $\pi/\Omega$ --- yakni laju Nyquist.
  \item Tunjukkan bahwa fungsi $x \mapsto
        \operatorname{sinc}(\Omega x - n\pi)$ dengan $n \in \Z$
        membentuk keluarga ortogonal di $L^2(\R)$ dengan norma
        tetap $\sqrt{\pi/\Omega}$, lalu turunkan identitas
        energinya $\norm f_2^2 =
        \frac\pi\Omega\sum_n\abs{f(n\pi/\Omega)}^2$.
  \item (Pelipatan spektrum) Tampilkan sebuah $g \in PW_{2\Omega}$ yang tak nol
        tetapi lenyap di setiap titik cuplikan $\frac{n\pi}\Omega$
        \emph{(tinjaulah $g(x) =
        \sin(\Omega x)\operatorname{sinc}(\Omega x)$ lalu
        periksalah pitanya)}: sehingga mencuplik di bawah laju Nyquist
        kehilangan keterangan --- sebab dua sinyal berbeda dapat berbagi
        semua cuplikannya: yakni gejala roda pedati stroboskopik,
        yang dimatematikakan.
  \item (Derajat kebebasan) Dengan memakai pertanyaan 14--15, benarkanlah
        kaidah keteknikannya: bahwa sinyal yang terbatas pita pada
        $\Omega$ dan energinya pada dasarnya terbawa oleh
        jendela waktu berpanjang $T$ diuraikan oleh
        kira-kira $\frac{\Omega T}\pi$ cuplikan real ---
        lalu persiskanlah ``pada dasarnya'' lewat identitas energinya
        dan ekornya $\sum_{\abs{n\pi/\Omega} >
        T/2}$.
  \item (Pemeriksaan kekonsistenan) Periksalah teorema pencuplikannya dengan
        tangan pada dua anggota $PW_\Omega$: (a) $f =
        \operatorname{sinc}(\Omega\,\cdot)$, yang cuplikannya
        $\delta_{n0}$; (b) sinyal pita sempit bertipe $f(x) = \cos(\omega x)
        \operatorname{sinc}(\varepsilon x)$ --- lebih tepatnya, tunjukkan bahwa untuk $f \in
        PW_{\Omega'}$ dengan $\Omega' < \Omega$, deret
        berlaju-$\Omega$ juga merekonstruksi $f$
        (sehingga pencuplikan berlebih tak berbahaya), lewat pembenaman $PW_
        {\Omega'} \subseteq PW_\Omega$.
\end{enumerate}

\medskip
\noindent\textbf{Bagian VI --- Ketakpastian, dua kali lagi.}
Ketaksamaan Heisenberg (\cref{exo:b3:fouriertransform:8})
membatasi \emph{seberapa} terpusat $f$ dan $\hat f$ secara bersama
dapat menjadi; dan berikut ini saudaranya yang serba-atau-tak-sama-sekali beserta
penjenuhannya yang persis.
\begin{enumerate}[resume]
  \item Misalkan $f \in L^1$ dengan $\operatorname{supp}f \subseteq
        \intcc{-A}A$. Tunjukkan bahwa $\hat f$ merupakan jumlah sebuah
        deret pangkat yang konvergen di mana-mana:
        \[
        \hat f(\xi) = \sum_{k\geq0}\frac{(-\iu\xi)^k}{k!}
        \,m_k, \qquad m_k = \int_{-A}^{A}x^kf(x)\,\dd x,
        \quad \abs{m_k} \leq A^k\norm f_1
        \]
        \emph{(uraikan $\eu^{-\iu\xi x}$ lalu benarkan
        penukarannya lewat kekonvergenan normalnya)}: sehingga transformasinya
        bersifat \emph{analitik real}, dengan jari-jari kekonvergenan tak
        berhingga di setiap titik.
  \item Turunkan dikotomi pendukungnya: bahwa fungsi analitik real
        yang lenyap pada sebuah selang terbuka tak kosong
        lenyap secara identik \emph{(sebab himpunan yang semua
        turunannya lenyap bersifat terbuka sekaligus tertutup --- kerjakanlah
        argumen Taylornya)}; lalu simpulkan bahwa tak ada $f$ tak nol yang
        $f$ maupun $\hat f$-nya berpendukung \omterm{def:b3:topology:compact}{kompak}, dan
        bahwa $PW_\Omega$ tak memuat fungsi berpendukung \omterm{def:b3:topology:compact}{kompak}
        yang tak nol --- sehingga sinyal terbatas pita berlangsung
        selamanya, dan sinyal terbatas waktu merembes ke semua
        frekuensinya.
  \item (Penjenuhan Heisenberg) Pada keluarga Gauss $f
        = \eu^{-ax^2}$, hitunglah kedua faktor pemusatannya
        lalu periksalah bahwa hasil kali ternormalkannya
        $\bigl(\int x^2\abs f^2\bigr)\bigl(\frac1{2\pi}\int
        \xi^2\abs{\hat f}^2\bigr)\big/\norm f_2^4$ sama dengan
        $\frac14$ untuk \emph{setiap} $a$ --- yakni keluarga kesamaan
        \cref{exo:b3:fouriertransform:8} yang sesungguhnya;
        lalu terangkan lewat argumen penskalaan mengapa hasil kalinya
        harus tetap di sepanjang keluarganya.
  \item Terangkan bacaan fisikanya (yakni \omterm{ex:b3:lebesgue:gamma}{kepadatan} posisi dan momentum
        sebuah keadaan kuantum; dan $\hbar$ pada penormalannya memberikan
        $\sigma_x\sigma_p \geq
        \frac\hbar2$), lalu hubungkanlah sepanjang bab ini:
        pemulusan seketika (Bagian II), ketakterbalikan (Bagian
        IV), pencuplikan (Bagian V), Heisenberg, dan dikotomi
        pendukungnya adalah lima ungkapan satu hukum --- yakni
        perilaku $\hat f$ di tak hingga mengatur apa yang boleh
        dilakukan $f$ di mana pun.
\end{enumerate}

\medskip
\noindent\textbf{Bagian VII --- Aljabar kernelnya, dan satu
contoh yang \omterm{def:b3:groups:derived}{terselesaikan}.}
\begin{enumerate}[resume]
  \item (Semigrup) Buktikan \emph{identitas
        Chapman--Kolmogorov} $g_t * g_s = g_{t+s}$ untuk $t, s > 0$ (lewat
        teorema konvolusinya dan keinjektifan transformasinya
        pada $L^1$), lalu turunkan $u(t + s) = g_s *
        u(t)$: sehingga berevolusi selama waktu $t + s$ adalah berevolusi selama $t$,
        lalu selama $s$. Lalu tajamkanlah pelesapan pertanyaan 8:
        dengan menulis $\norm{u(t)}_2^2 = \frac1{2\pi}\int
        \eu^{-2t\xi^2}\abs{\hat f(\xi)}^2\dd\xi$, tunjukkan lewat
        Cauchy--Schwarz bahwa
        \[
        t \longmapsto \ln\,\norm{u(t)}_2
        \quad\text{bersifat cembung pada } \intoo0{+\infty} :
        \]
        sehingga energi $L^2$-nya bukan sekadar menurun, melainkan
        menurun secara log-cembung.
  \item (Ke mana panasnya pergi) Misalkan $f \geq 0$ dengan $f \in L^1$
        dan $\int x^2f(x)\dd x < \infty$. Tunjukkan bahwa untuk setiap
        $t > 0$
        \[
        \int_\R x\,u(t, x)\,\dd x = \int_\R x f(x)\,\dd x,
        \qquad
        \int_\R x^2u(t, x)\,\dd x
        = \int_\R x^2f(x)\,\dd x + 2t\int_\R f :
        \]
        sehingga pusat panasnya tak pernah bergerak, dan ragamnya
        tumbuh \emph{secara linear} terhadap waktu --- yakni penskalaan
        difusif $x \sim \sqrt{2t}$, yang akan dibaca ulang saat
        gerak Brown muncul di \cref{ch:b3:probability}.
        \emph{(Hitunglah dua momen pertama $g_t$ lalu pakailah
        Tonelli pada konvolusinya.)}
  \item (Fungsi Gaussnya, yang diselesaikan sampai tuntas) Untuk $f(x) =
        \eu^{-x^2}$, tegakkanlah bentuk tertutupnya
        \[
        u(t, x) = \frac1{\sqrt{1 + 4t}}\,
        \exp\Bigl(-\frac{x^2}{1 + 4t}\Bigr),
        \]
        lalu periksalah padanya, dengan tangan: persamaan $\partial_tu
        = \partial^2_{xx}u$; kekekalan $\int u(t) =
        \sqrt\pi$; hukum pelesapan $\norm{u(t)}_2 =
        (\pi/2)^{1/4}(1 + 4t)^{-1/4}$ (bandingkan peluruhan $t^{-1/4}$-nya
        dengan peluruhan norma-sup $t^{-1/2}$ pada pertanyaan
        8); dan pertumbuhan ragam yang persis pada pertanyaan 24. Pada
        $t = 6$: puncaknya telah turun menjadi $\frac15$ tinggi
        awalnya sedangkan profilnya lima kali lebih lebar
        --- yakni panas yang sama, yang tersebar.

\end{enumerate}
\end{problem}
